% !TEX spellcheck = en_US
\documentclass[aspectratio=169,10pt,compress]{beamer}
\usepackage{../../cis000_latex/lib/hblibPresentation}
\usepackage{../../cis000_latex/lib/hblibCodeoutput}
\usepackage{../../cis000_latex/lib/hblibDatastructures}
\usepackage[T1]{fontenc}
\usepackage{lmodern}
\usepackage{booktabs}

\lstset{columns=fixed, breaklines=true, breakatwhitespace=false}

\title{Linked Stacks, Queues, and Deques}
\subtitle{Concepts, Implementation, and Applications}
\author[Bingol]{Ruhan Alpaydin}
\institute[FCIS]{Faculty of Computer and Information Sciences\\Yeditepe University}
\date{{\tiny \hbTimeStamp}}

\begin{document}

\begin{frame}[t,fragile]\titlepage\end{frame}

\begin{frame}[t,fragile]\frametitle{Outline}
  \begin{columns}[T]
    \begin{column}{0.48\textwidth}
      \tableofcontents[sections={1-2},hideallsubsections]
    \end{column}
    \begin{column}{0.48\textwidth}
      \tableofcontents[sections={3},hideallsubsections]
    \end{column}
  \end{columns}
\end{frame}

% =====================================================================
% PART 1: CONCEPTS (from concepts.tex)
% =====================================================================
\hSection[Nodes]{Concepts: Nodes vs Arrays}

\begin{frame}[t,fragile]\frametitle{Two ways to build a container}
  \begin{columns}[T]
    \begin{column}{0.48\textwidth}
      \begin{block}{Array-backed}
        \small
        Elements in one contiguous block.\\
        Fast random access: $O(1)$.\\
        Insert/delete at ends can require shifting or resizing.
      \end{block}
    \end{column}
    \begin{column}{0.48\textwidth}
      \begin{block}{Node-backed (Linked)}
        \small
        Each element in its own node with a reference.\\
        Random access is $O(n)$ (walk).\\
        Insert/delete at known ends is always $O(1)$.
      \end{block}
    \end{column}
  \end{columns}
  \vspace{0.4em}
  \small
  \hRemark{We trade random access for guaranteed $O(1)$ end-operations.}
\end{frame}

\begin{frame}[t,fragile]\frametitle{The three ADTs}
  \begin{center}\small
  \begin{tabular}{@{}lll@{}}
    \toprule
    \textbf{ADT} & \textbf{Rule} & \textbf{Core Operations} \\
    \midrule
    Stack & LIFO (Last In, First Out) & \hCode{push}, \hCode{pop} \\
    Queue & FIFO (First In, First Out) & \hCode{enqueue}, \hCode{dequeue} \\
    Deque & Double-Ended Queue & \hCode{addFirst}, \hCode{addLast}, \hCode{removeFirst}, \hCode{removeLast} \\
    \bottomrule
  \end{tabular}
  \end{center}
\end{frame}

\begin{frame}[t,fragile]\frametitle{Stack: LIFO}
  \begin{columns}[T]
    \begin{column}{0.48\textwidth}
      \begin{itemize}\small
        \item \hCode{push(x)}: add a new node at the \textbf{top}.
        \item \hCode{pop()}: remove the top node.
        \item Always $O(1)$, no resizing.
      \end{itemize}
    \end{column}
    \begin{column}{0.48\textwidth}
      \begin{center}
      \begin{tikzpicture}[arr/.style={-{Latex[length=2mm]}, thick}]
        \node at (-1,0) {\hCode{top}};
        \node[draw, minimum width=1.5cm, minimum height=0.6cm] (a) at (0.5,0) {10};
        \node[draw, minimum width=1.5cm, minimum height=0.6cm] (b) at (2.5,0) {20};
        \node at (4.0,0) {\hCode{null}};
        \draw[arr] (-0.3,0) -- (a.west);
        \draw[arr] (a.east) -- (b.west);
        \draw[arr] (b.east) -- (3.8,0);
      \end{tikzpicture}
      \end{center}
      \small \hCode{pop} returns \textbf{10}.
    \end{column}
  \end{columns}
\end{frame}

\begin{frame}[t,fragile]\frametitle{Queue: FIFO}
  \begin{columns}[T]
    \begin{column}{0.48\textwidth}
      \begin{itemize}\small
        \item \hCode{enqueue(x)}: add a new node at the \textbf{tail}.
        \item \hCode{dequeue()}: remove the \textbf{head} node.
        \item Always $O(1)$, provided we keep a \hCode{tail} reference.
      \end{itemize}
    \end{column}
    \begin{column}{0.48\textwidth}
      \begin{center}
      \begin{tikzpicture}[arr/.style={-{Latex[length=2mm]}, thick}]
        \node at (-1,0) {\hCode{head}};
        \node[draw, minimum width=1.5cm, minimum height=0.6cm] (a) at (0.5,0) {10};
        \node[draw, minimum width=1.5cm, minimum height=0.6cm] (b) at (2.5,0) {20};
        \node at (4.0,0) {\hCode{null}};
        \node at (2.5,-1) {\hCode{tail}};
        \draw[arr] (-0.3,0) -- (a.west);
        \draw[arr] (a.east) -- (b.west);
        \draw[arr] (b.east) -- (3.8,0);
        \draw[arr] (2.5,-0.8) -- (b.south);
      \end{tikzpicture}
      \end{center}
      \small \hCode{dequeue} returns \textbf{10}.
    \end{column}
  \end{columns}
\end{frame}

\begin{frame}[t,fragile]\frametitle{Deque: double-ended}
  \begin{columns}[T]
    \begin{column}{0.48\textwidth}
      \begin{itemize}\small
        \item \hCode{addFirst(x)}, \hCode{addLast(x)}.
        \item \hCode{removeFirst()}, \hCode{removeLast()}.
        \item To do \hCode{removeLast()} in $O(1)$, the nodes must be
              \textbf{doubly linked} (each node needs a \hCode{prev}).
      \end{itemize}
    \end{column}
    \begin{column}{0.48\textwidth}
      \begin{center}
      \begin{tikzpicture}[arr/.style={-{Latex[length=2mm]}, thick}]
        \node at (-1,0) {\hCode{head}};
        \node[draw, minimum width=1.5cm, minimum height=0.6cm] (a) at (0.5,0) {10};
        \node[draw, minimum width=1.5cm, minimum height=0.6cm] (b) at (2.5,0) {20};
        \node at (4.0,0) {\hCode{null}};
        \node at (2.5,-1) {\hCode{tail}};
        \draw[arr] (-0.3,0) -- (a.west);
        \draw[arr] (a.east) -- (b.west);
        \draw[arr] (b.east) -- (3.8,0);
        \draw[arr] (2.5,-0.8) -- (b.south);
        \draw[arr] (b.south west) -- (a.south east);
      \end{tikzpicture}
      \end{center}
      \small \hCode{removeLast} returns \textbf{20}.
    \end{column}
  \end{columns}
\end{frame}

\begin{frame}[t,fragile]\frametitle{The Java API}
  \begin{center}\small
  \begin{tabular}{@{}lll@{}}
    \toprule
    \textbf{Interface} & \textbf{Implements} & \textbf{Note} \\
    \midrule
    \hCode{Deque<E>} & \hCode{ArrayDeque}, \hCode{LinkedList} & Covers Stack and Queue \\
    \hCode{Queue<E>} & \hCode{LinkedList}, \hCode{ArrayDeque} & FIFO operations \\
    \emph{(none)} & \hCode{java.util.Stack} & Legacy, avoid it \\
    \bottomrule
  \end{tabular}
  \end{center}
  \vspace{0.3em}
  \begin{itemize}\small
    \item \hCode{ArrayDeque} is the modern default: fast, memory-efficient.
    \item \hCode{LinkedList} is a doubly linked list. Use it when you
          specifically want a linked structure or also need \hCode{List}.
  \end{itemize}
\end{frame}

% =====================================================================
% PART 2: FRESHMAN (from freshman.tex)
% =====================================================================
\hSection[Impl]{Implementation}

\begin{frame}[t,fragile]\frametitle{The \texttt{Node} Class}
  \begin{columns}[T]
    \begin{column}{0.5\textwidth}
      \lstcode[basicstyle=\ttfamily\scriptsize]{codeoutput/Node.java}{Java}
    \end{column}
    \begin{column}{0.46\textwidth}
      \begin{itemize}\small
        \item One data field: \hCode{int data}.
        \item One reference: \hCode{Node next}.
        \item The constructor leaves \hCode{next} as \hCode{null}.
      \end{itemize}
    \end{column}
  \end{columns}
\end{frame}

\begin{frame}[t,fragile]\frametitle{\texttt{LinkedStack}: push}
  \begin{columns}[T]
    \begin{column}{0.52\textwidth}
      \lstcode[basicstyle=\ttfamily\scriptsize, firstline=4, lastline=8, firstnumber=4]{codeoutput/LinkedStack.java}{Java}
    \end{column}
    \begin{column}{0.44\textwidth}
      \begin{block}{Two steps}
        \small
        1. new node's \hCode{next} $\gets$ current \hCode{top}.\\
        2. \hCode{top} $\gets$ new node.
      \end{block}
      \begin{block}{Cost}
        \small
        $O(1)$ time. No resize.
      \end{block}
    \end{column}
  \end{columns}
\end{frame}

\begin{frame}[t,fragile]\frametitle{\texttt{LinkedStack}: pop}
%  \begin{columns}[T]
 %   \begin{column}{0.52\textwidth}
      \lstcode[basicstyle=\ttfamily\scriptsize, firstline=10, lastline=15, firstnumber=10]{codeoutput/LinkedStack.java}{Java}
 %   \end{column}
 %   \begin{column}{0.44\textwidth}
      \begin{block}{Three steps}
        \small
        1. save data from \hCode{top}.\\
        2. move \hCode{top} to \hCode{top.next}.\\
        3. return data.
      \end{block}
      \begin{block}{Underflow}
        \small
        \hCode{top == null}: throw exception.
      \end{block}
 %   \end{column}
  %\end{columns}
\end{frame}

\begin{frame}[t,fragile]\frametitle{\texttt{LinkedStack}: run}
  \begin{columns}[T]
    \begin{column}{0.52\textwidth}
      \lstcode[basicstyle=\ttfamily\tiny]{codeoutput/LinkedStack.java}{Java}
    \end{column}
    \begin{column}{0.44\textwidth}
      \lstterminal[basicstyle=\ttfamily\scriptsize]{codeoutput/LinkedStack.txt}
      \small
      LIFO order: 30, 20, 10.
    \end{column}
  \end{columns}
\end{frame}

\begin{frame}[t,fragile]\frametitle{\texttt{LinkedQueue}: enqueue}
%  \begin{columns}[T]
%    \begin{column}{0.52\textwidth}
      \lstcode[basicstyle=\ttfamily\scriptsize, firstline=5, lastline=10, firstnumber=5]{codeoutput/LinkedQueue.java}{Java}
%    \end{column}
%    \begin{column}{0.44\textwidth}
      \begin{block}{Append at tail}
        \small
        If empty, \hCode{head = tail = n}.\\
        Else \hCode{tail.next = n}, then \hCode{tail = n}.
      \end{block}
      \begin{block}{Cost}
        \small
        $O(1)$ because we keep a \hCode{tail} reference.
      \end{block}
%    \end{column}
%  \end{columns}
\end{frame}

\begin{frame}[t,fragile]\frametitle{\texttt{LinkedQueue}: dequeue}
 % \begin{columns}[T]
 %   \begin{column}{0.52\textwidth}
      \lstcode[basicstyle=\ttfamily\scriptsize, firstline=11, lastline=17, firstnumber=11]{codeoutput/LinkedQueue.java}{Java}
%    \end{column}
 %   \begin{column}{0.44\textwidth}
      \begin{block}{Remove at head}
        \small
        Save \hCode{head.data}.\\
        Advance \hCode{head = head.next}.\\
        If \hCode{head == null}, set \hCode{tail = null}.
      \end{block}
      \begin{block}{Critical}
        \small
        Forgetting \hCode{tail = null} leaves a dangling tail!
      \end{block}
 %   \end{column}
%  \end{columns}
\end{frame}

\begin{frame}[t,fragile]\frametitle{\texttt{LinkedQueue}: run}
  \begin{columns}[T]
    \begin{column}{0.52\textwidth}
      \lstcode[basicstyle=\ttfamily\tiny]{codeoutput/LinkedQueue.java}{Java}
    \end{column}
    \begin{column}{0.44\textwidth}
      \lstterminal[basicstyle=\ttfamily\scriptsize]{codeoutput/LinkedQueue.txt}
      \small
      FIFO order: 10, 20, 30.
    \end{column}
  \end{columns}
\end{frame}

\begin{frame}[t,fragile]\frametitle{\texttt{LinkedDeque}: \texttt{addFirst} and \texttt{addLast}}
  \begin{columns}[T]
    \begin{column}{0.52\textwidth}
      \lstcode[basicstyle=\ttfamily\scriptsize, firstline=5, lastline=18, firstnumber=5]{codeoutput/LinkedDeque.java}{Java}
    \end{column}
    \begin{column}{0.44\textwidth}
      \begin{itemize}\small
        \item \hCode{addFirst}: insert before \hCode{head}.
        \item \hCode{addLast}: append after \hCode{tail}.
        \item Both $O(1)$.
        \item To do \hCode{removeLast} in $O(1)$ we would need
              a doubly linked node.
      \end{itemize}
    \end{column}
  \end{columns}
\end{frame}

\begin{frame}[t,fragile]\frametitle{\texttt{LinkedDeque}: run}
  \begin{columns}[T]
    \begin{column}{0.56\textwidth}
      \lstcode[basicstyle=\ttfamily\tiny]{codeoutput/LinkedDeque.java}{Java}
    \end{column}
    \begin{column}{0.42\textwidth}
      \lstterminal[basicstyle=\ttfamily\scriptsize]{codeoutput/LinkedDeque.txt}
      \small
      We added 20, then 10 at front, then 30 at back.\\
      Removing from front yields 10, then 20.
    \end{column}
  \end{columns}
\end{frame}

% =====================================================================
% PART 3: EXAMPLES (from examples.tex)
% =====================================================================
\hSection[Apps]{Applications}

\begin{frame}[t,fragile]\frametitle{Valid Brackets}
  \begin{columns}[T]
    \begin{column}{0.54\textwidth}
      \lstcode[basicstyle=\ttfamily\tiny]{codeoutput/Brackets.java}{Java}
    \end{column}
    \begin{column}{0.42\textwidth}
      \lstterminal[basicstyle=\ttfamily\scriptsize]{codeoutput/Brackets.txt}
      \small
      Push opening brackets. Pop and match on closing.
      If the stack is empty at the end, it is valid.
    \end{column}
  \end{columns}
\end{frame}

\begin{frame}[t,fragile]\frametitle{Palindrome}
  \begin{columns}[T]
    \begin{column}{0.54\textwidth}
      \lstcode[basicstyle=\ttfamily\tiny]{codeoutput/Palindrome.java}{Java}
    \end{column}
    \begin{column}{0.42\textwidth}
      \lstterminal[basicstyle=\ttfamily\scriptsize]{codeoutput/Palindrome.txt}
      \small
      Load characters into a Deque. Compare and remove from both ends.
      If it finishes without a mismatch, it is a palindrome.
    \end{column}
  \end{columns}
\end{frame}

\begin{frame}[t,fragile]
  \begin{center}\vspace{4cm}\hRemark{\Huge Questions?}\end{center}
\end{frame}

\end{document}
